\section{The variance for a single modulus}\label{sec:single}

\begin{proposition}[Single-modulus variance]\label{prop:single}
For every $\eps>0$ there is $c_7(\eps)$ such that for all real $x\ge2$ and all integers $1\le q\le x$:
\begin{enumerate}[label=\textup{(\roman*)}]
\item $|V(x;q)-\MT(x;q)|\le c_7(\eps)\bigl(x^{5/3+\eps}/q+\sqrt x\bigr)$;
\item $|V(x;q)-h(q)\sqrt{xq}|\le c_7(\eps)\bigl(x^{1/3}q^{2/3+\eps}+x^{5/3+\eps}/q+\sqrt x\bigr)$.
\end{enumerate}
\end{proposition}

\begin{proof}
(i) By \cref{prop:identity}, which holds for all $x\ge1$ and $q\in\N$,
\[
V(x;q)-\MT(x;q)=2\sum_{1\le h\le x/q}D_x(hq)+\Bigl(Q(x)-\frac{6x}{\pi^2}\Bigr)-2x\sum_{l=1}^qg(q,l)E(x;q,l).
\]
For $1\le h\le x/q$ the integer $m=hq$ satisfies $1\le m\le x$. Since $x\ge2$, \cref{prop:pairs} applies to each of the at most $x/q$ terms. (This is the only place where $x\ge2$ rather than $x\ge1$ is needed.) It gives
\[
\Bigl|2\sum_{1\le h\le x/q}D_x(hq)\Bigr|\le2\cdot\frac xq\cdot c_3(\eps)x^{2/3+\eps}=2c_3(\eps)\frac{x^{5/3+\eps}}q .
\]
\Cref{prop:trivial}(i) gives $|Q(x)-6x/\pi^2|\le c_0\sqrt x$. \Cref{prop:trivial}(ii) gives
\[
\Bigl|2x\sum_lg(q,l)E(x;q,l)\Bigr|\le2c_0\frac{x^{3/2}}q\le2c_0\frac{x^{5/3+\eps}}q,
\]
since $x\ge1$. By the triangle inequality, (i) holds with $c_7(\eps)\ge2c_3(\eps)+2c_0$.

(ii) \Cref{prop:MT} gives $|\MT(x;q)-h(q)\sqrt{xq}|\le c_6(\eps)x^{1/3}q^{2/3+\eps}$. Combining this with (i) by the triangle inequality proves (ii) with
\[
c_7(\eps)=\max\bigl\{2c_3(\eps)+2c_0,\ c_6(\eps)\bigr\},
\]
and (i) remains true with this larger constant.
\end{proof}
